% part: intuitionistic-logic
% chapter: soundness-completeness
% section: canonical-model

\documentclass[../../../include/open-logic-section]{subfiles}

\begin{document}

\olfileid{int}{sc}{mod}

\olsection{కానానికల్ నమూనా}

మన నమూనాలోని లోకాలు ప్రకృతిసంఖ్యల
పరిమిత క్రమాలు~$\sigma$; అంటే
$\sigma \in \Nat^*$. $\Nat^*$ను
కింది విధంగా ఆగమన పద్ధతిలో
నిర్వచిస్తామని గమనించండి:
\begin{enumerate}
\item $\emptyseq \in \Nat^*$.
\item $\sigma \in \Nat^*$ మరియు $n \in \Nat$
  అయితే $\sigma.n \in
  \Nat^*$ (ఇక్కడ $\sigma.n$ అనేది
  $\sigma \concat \tuple{n}$; $\sigma \concat \sigma'$
  అనేది $\sigma$, $\sigma'$ల అనుసంధానం).
\item వీటికి మించి ఏదీ $\Nat^*$లో లేదు.
\end{enumerate}
కాబట్టి $\Nat^*$ను ఆగమన నిర్వచనాల కోసం
వాడవచ్చు.

\tetoken{సూత్రాల}{!!{formula}s} అన్ని జతలను
$\tuple{!B_1, !C_1}$, $\tuple{!B_2, !C_2}$,
\dots, అని వరుసపెడదాం. \tetoken{సూత్రాల}{!!{formula}s}
సమితి~$\Delta$ ఇచ్చినప్పుడు
$\Delta(\sigma)$ను ఇలా ఆగమన పద్ధతిలో
నిర్వచిద్దాం:
\begin{enumerate}
\item $\Delta(\emptyseq) = \Delta$
\item $\Delta(\sigma.n) = {}$
  \[
  \begin{cases}
    (\Delta(\sigma) \cup \{!B_n\})^* &
    \text{$\Delta(\sigma) \cup \{!B_n\} \Proves/ !C_n$ అయితే} \\
    \Delta(\sigma) & \text{లేకపోతే}
  \end{cases}
  \]
\end{enumerate}
ఇక్కడ $(\Delta(\sigma) \cup \{!B_n\})^*$
అంటే, $\Delta(\sigma) \cup \{!B_n\}$ సమితికి,
\tetoken{సూత్రం}{!!{formula}}~$!C_n$కు
\olref[lin]{lem:lindenbaum}ను వర్తింపజేస్తే
లభించే ప్రధాన \tetoken{సూత్రాల}{!!{formula}s}
సమితి. ఈ నిర్వచనం ప్రకారం
$\Delta(\sigma) \cup \{!B_n\} \Proves/ !C_n$
అయితే $\Delta(\sigma.n)
\Proves !B_n$ మరియు
$\Delta(\sigma.n) \Proves/ !C_n$ అని
గమనించండి. ఏ~$n$కైనా
$\Delta(\sigma) \subseteq \Delta(\sigma.n)$
అని కూడా గమనించండి. $\Delta$
ప్రధానమైతే, అన్ని~$\sigma$లకు
$\Delta(\sigma)$ ప్రధానమే.

\begin{defn}\ollabel{defn:canonical-model}
  $\Delta$ ప్రధానమని అనుకుందాం. అప్పుడు
  $\Delta$కు \emph{కానానికల్ నమూనా}
  $\mModel{M(\Delta)}$ ఇలా నిర్వచించబడుతుంది:
  \begin{enumerate}
  \item $W = \Nat^*$; ప్రకృతిసంఖ్యల పరిమిత
    క్రమాల సమితి.
  \item $R$ అనేది $R\sigma\sigma'$
    అయితే, అప్పుడే $\sigma$ అనేది~$\sigma'$
    యొక్క ప్రారంభ భాగం అయ్యే పాక్షిక క్రమం;
    అంటే ఏదో క్రమం~$\sigma''$కు
    $\sigma' =
    \sigma \concat \sigma''$.
  \item $V(p) = \Setabs{\sigma}{p \in \Delta(\sigma)}$.
  \end{enumerate}
\end{defn}

$R$ నిజంగానే పాక్షిక క్రమమని సులభంగా
తనిఖీ చేయవచ్చు. అలాగే~$V$కు ఏకదిశ
పెరుగుదల షరతు నెరవేరుతుంది.
$\Delta(\sigma)
\subseteq \Delta(\sigma.n)$ కాబట్టి,
$R\sigma\sigma'$ అయినప్పుడల్లా
$\Delta(\sigma)
\subseteq \Delta(\sigma')$ వస్తుంది:
ఇక్కడ~$\sigma$ను స్థిర మూల క్రమంగా తీసుకొని,
జోడించిన భాగపు పొడవుపై ఆగమన పద్ధతిలో
నిరూపించవచ్చు.
\sourcecorrection{OLTEINTCAN-001}{మూల చివరి వాక్యం $\sigma$పై ఆగమనం అంటుంది; కానీ ఒకే మూల $\sigma$కు చెందిన అనేక విస్తరణలలో చేర్పు నిరూపణ జోడించిన పరిమిత భాగపు పొడవుపై సాగుతుంది.}

\end{document}
